% \iffalse meta-comment
%
% hit-option-engine.dtx 是各个类共用的选项声明与判断助手
%
% \fi
%
% \section{选项机制（engine）}
%
% \changes{v3.2a}{2026/08/11}{声明选项的助手与判断函数抽出来共用；具名判断由
%   \cs{hit@define@condition} 生成（原先每条十来行样板）；加
%   \cs{hit@define@boolean@option@as}；基础类加载之后再遇到不认识的键报警告；只是
%   expl3 别名的几个判断与一处没调用的宏删掉}
%
% 提供类选项声明和条件判断函数。
%
%    \begin{macrocode}
%<*shared-option-engine>
%    \end{macrocode}
%
% 声明选项
%
% \cs{hit@define@boolean@option}\marg{名}\marg{默认} 声明一个布尔选项，标志位是
% \cs{g\_\_hit\_}\meta{名}\cs{\_bool}；\cs{hit@define@string@option}\marg{名} 声明一个存字符串的选项，
% 存进 \cs{hit@}\meta{名}。键名带连字符时不能直接拿来拼控制序列，
% \cs{hit@define@string@option@as}\marg{键名}\marg{宏名} 把两者分开。
%
% 布尔选项也有一个分开写的版本
% \cs{hit@define@boolean@option@as}\marg{键名}\marg{标志位名}\marg{默认}。
% 改过名的选项要用它：\option{arialtoc} 这个键名留在兼容期，标志位却该按新名
% \cs{tableofcontents}\oarg{latin-sans} 叫，两套名字指的是同一个开关。
%
% \cs{hit@define@choice@option}\marg{键名}\marg{前缀}\marg{取值表}\marg{默认} 声明一组互斥取值，
% 每个取值一个标志位，选中的置真、其余置假，取值不在表里就报错。
%    \begin{macrocode}
%    \end{macrocode}
%
% 基础类加载前暂存未知选项以便转发，加载后遇到未知键则警告。
%    \begin{macrocode}
\bool_new:N \g__hit_base_class_loaded_bool
\clist_new:N \g__hit_base_class_options_clist
\cs_generate_variant:Nn \msg_warning:nnn { nnV }
\cs_new_protected:Npn \hit@pass@or@warn #1
  {
    \bool_if:NTF \g__hit_base_class_loaded_bool
      { \msg_warning:nnV { hithesis } { unknown-key } \l_keys_key_str }
      {
        \tl_if_empty:nTF {#1}
          { \clist_gput_right:NV \g__hit_base_class_options_clist \l_keys_key_str }
          {
            \clist_gput_right:Nx \g__hit_base_class_options_clist
              { \l_keys_key_str = \exp_not:n {#1} }
          }
      }
  }
\cs_new_protected:Npn \hit@define@boolean@option #1#2
  { \hit@define@boolean@option@as {#1} {#1} {#2} }
\cs_new_protected:Npn \hit@define@boolean@option@as #1#2#3
  {
    \bool_new:c { g__hit_#2_bool }
    \str_if_eq:nnT {#3} { true } { \bool_gset_true:c { g__hit_#2_bool } }
    \keys_define:nn { hit / class }
      { #1 .bool_gset:c = { g__hit_#2_bool } , #1 .default:n = { true } }
  }
\cs_new_protected:Npn \hit@define@string@option #1 { \hit@define@string@option@as {#1} {#1} }
\cs_new_protected:Npn \hit@define@string@option@as #1#2
  {
    \tl_new:c { hit@#2 }
    \keys_define:nn { hit / class } { #1 .tl_gset:c = { hit@#2 } }
  }
\cs_new_protected:Npn \hit@define@choice@option #1#2#3#4
  {
    \clist_map_inline:nn {#3} { \bool_new:c { g__hit_#2##1_bool } }
    \tl_if_empty:nF {#4} { \bool_gset_true:c { g__hit_#2#4_bool } }
    \keys_define:nn { hit / class }
      {
        #1 .code:n =
          {
            \clist_map_inline:nn {#3}
              { \bool_gset_false:c { g__hit_#2####1_bool } }
            \bool_if_exist:cTF { g__hit_#2##1_bool }
              { \bool_gset_true:c { g__hit_#2##1_bool } }
              { \hit@bad@value:nnn {#1} {##1} {#3} }
          } ,
      }
  }
%    \end{macrocode}
%
% 读标志位的判断函数
%
% 这些条件命令可在普通 LaTeX 语法中调用，参数内的空格按调用处的 catcode 保留。
%
% 名字末尾的 \texttt{@TF} 表示真假两个分支都要给，\texttt{@T} 表示只有真分支。
% \cs{hit@if@option@TF}\marg{名}\marg{真}\marg{假} 读一个布尔选项的标志位，
% \cs{hit@if@both@options@TF}\marg{名}\marg{名}\marg{真}\marg{假} 要两个同时成立，
%
%    \begin{macrocode}
\cs_new:Npn \hit@if@option@TF #1#2#3 { \bool_if:cTF { g__hit_#1_bool } {#2} {#3} }
\cs_new:Npn \hit@if@option@T #1#2 { \bool_if:cT { g__hit_#1_bool } {#2} }
\cs_new:Npn \hit@if@both@options@TF #1#2#3#4
  {
    \bool_lazy_and:nnTF
      { \bool_if_p:c { g__hit_#1_bool } } { \bool_if_p:c { g__hit_#2_bool } } {#3} {#4}
  }
%    \end{macrocode}
%
% \cs{hit@define@condition}\marg{名}\marg{布尔表达式} 造一对判断：
% \cs{hit@if@}\meta{名} 只有真分支，\cs{hit@if@}\meta{名}\texttt{@TF} 两个分支都有。
% 表达式用 \pkg{expl3} 的布尔写法，\texttt{\&\&}、\texttt{||}、\texttt{!} 与括号都能用。
%
% 名字里的词用 \texttt{@} 分隔、一个都不缩写，虚词也留着（\texttt{harbin@master@or@doctor}）。
% 控制序列里不能用连字符，\texttt{@} 是这里的等价物。
%
% 具体条件在选项模块中声明。
%    \begin{macrocode}
\cs_new_protected:Npn \hit@define@condition #1#2
  {
    \cs_new:cpn { hit@if@#1 } ##1 { \bool_if:nT {#2} {##1} }
    \cs_new:cpn { hit@if@#1 @TF } ##1##2 { \bool_if:nTF {#2} {##1} {##2} }
  }
%    \end{macrocode}
%
% \option{style/glue} 开时给行距一个伸缩量，关时是固定值。只剩普通报告的正文行距在用，终稿
% 的伸缩量由 \option{Normal} 样式的行距区间给。
%    \begin{macrocode}
\cs_new:Npn \hit@glue@skip #1#2
  { \bool_if:NTF \g__hit_glue_bool { #1~plus~#2~minus~0bp } { #1 } }
%    \end{macrocode}
%
%    \begin{macrocode}
%</shared-option-engine>
%    \end{macrocode}
%
% \Finale
